\documentclass[10pt,a4paper]{amsart}
\usepackage{amsmath,amssymb,amsthm,graphicx}
\usepackage{verbatim}
\newenvironment{reference}{\comment}{\endcomment}
%\newenvironment{reference}{}{}
\newenvironment{slogan}{\comment}{\endcomment}
\newenvironment{history}{\comment}{\endcomment}

% For commutative diagrams we use Xy-pic
\usepackage[all]{xy}

% We use 2cell for 2-commutative diagrams.
\xyoption{2cell}
% Theorem environments.
%
\theoremstyle{plain}
\newtheorem{theorem}[subsection]{Theorem}
\newtheorem{proposition}[subsection]{Proposition}
\newtheorem{lemma}[subsection]{Lemma}

\theoremstyle{definition}
\newtheorem{definition}[subsection]{Definition}
\newtheorem{example}[subsection]{Example}
\newtheorem{exercise}[subsection]{Exercise}
\newtheorem{situation}[subsection]{Situation}

\theoremstyle{remark}
\newtheorem{remark}[subsection]{Remark}
\newtheorem{remarks}[subsection]{Remarks}

\numberwithin{equation}{subsection}

% Macros
%
\def\lim{\mathop{\mathrm{lim}}\nolimits}
\def\colim{\mathop{\mathrm{colim}}\nolimits}
\def\Spec{\mathop{\mathrm{Spec}}}
\def\Hom{\mathop{\mathrm{Hom}}\nolimits}
\def\Ext{\mathop{\mathrm{Ext}}\nolimits}
\def\SheafHom{\mathop{\mathcal{H}\!\mathit{om}}\nolimits}
\def\SheafExt{\mathop{\mathcal{E}\!\mathit{xt}}\nolimits}
\def\Sch{\mathit{Sch}}
\def\Mor{\mathop{\mathrm{Mor}}\nolimits}
\def\Ob{\mathop{\mathrm{Ob}}\nolimits}
\def\Sh{\mathop{\mathit{Sh}}\nolimits}
\def\NL{\mathop{N\!L}\nolimits}
\def\CH{\mathop{\mathrm{CH}}\nolimits}
\def\proetale{{pro\text{-}\acute{e}tale}}
\def\etale{{\acute{e}tale}}
\def\QCoh{\mathit{QCoh}}
\def\Ker{\mathop{\mathrm{Ker}}}
\def\Im{\mathop{\mathrm{Im}}}
\def\Coker{\mathop{\mathrm{Coker}}}
\def\Coim{\mathop{\mathrm{Coim}}}

% Boxtimes
%
\DeclareMathSymbol{\boxtimes}{\mathbin}{AMSa}{"02}

%
% Macros for moduli stacks/spaces
%
\def\QCohstack{\mathcal{QC}\!\mathit{oh}}
\def\Cohstack{\mathcal{C}\!\mathit{oh}}
\def\Spacesstack{\mathcal{S}\!\mathit{paces}}
\def\Quotfunctor{\mathrm{Quot}}
\def\Hilbfunctor{\mathrm{Hilb}}
\def\Curvesstack{\mathcal{C}\!\mathit{urves}}
\def\Polarizedstack{\mathcal{P}\!\mathit{olarized}}
\def\Complexesstack{\mathcal{C}\!\mathit{omplexes}}
% \Pic is the operator that assigns to X its picard group, usage \Pic(X)
% \Picardstack_{X/B} denotes the Picard stack of X over B
% \Picardfunctor_{X/B} denotes the Picard functor of X over B
\def\Pic{\mathop{\mathrm{Pic}}\nolimits}
\def\Picardstack{\mathcal{P}\!\mathit{ic}}
\def\Picardfunctor{\mathrm{Pic}}
\def\Deformationcategory{\mathcal{D}\!\mathit{ef}}
\usepackage{cleveref}
\crefname{originalreference}{original reference}{original references}
\Crefname{originalreference}{Original reference}{Original references}
\makeatletter
\newlabel{situation-setup}{{02QL}{source}{}{sourcepointer-1}{}}
\newlabel{situation-setup@cref}{{[originalreference][0][]02QL}{[1][0][]source}{}{}{}}
\makeatother
\crefname{originallemma}{original lemma}{original lemmas}
\Crefname{originallemma}{Original lemma}{Original lemmas}
\makeatletter
\newlabel{lemma-flat-pullback-bivariant}{{0B78}{source}{}{sourcepointer-2}{}}
\newlabel{lemma-flat-pullback-bivariant@cref}{{[originallemma][0][]0B78}{[1][0][]source}{}{}{}}
\makeatother
\crefname{originallemma}{original lemma}{original lemmas}
\Crefname{originallemma}{Original lemma}{Original lemmas}
\makeatletter
\newlabel{morphisms-lemma-base-change-smooth}{{01VB}{source}{}{sourcepointer-3}{}}
\newlabel{morphisms-lemma-base-change-smooth@cref}{{[originallemma][0][]01VB}{[1][0][]source}{}{}{}}
\makeatother
\crefname{originallemma}{original lemma}{original lemmas}
\Crefname{originallemma}{Original lemma}{Original lemmas}
\makeatletter
\newlabel{morphisms-lemma-base-change-syntomic}{{01UI}{source}{}{sourcepointer-4}{}}
\newlabel{morphisms-lemma-base-change-syntomic@cref}{{[originallemma][0][]01UI}{[1][0][]source}{}{}{}}
\makeatother
\crefname{originallemma}{original lemma}{original lemmas}
\Crefname{originallemma}{Original lemma}{Original lemmas}
\makeatletter
\newlabel{more-morphisms-lemma-flat-lci}{{069K}{source}{}{sourcepointer-5}{}}
\newlabel{more-morphisms-lemma-flat-lci@cref}{{[originallemma][0][]069K}{[1][0][]source}{}{}{}}
\makeatother
\crefname{originallemma}{original lemma}{original lemmas}
\Crefname{originallemma}{Original lemma}{Original lemmas}
\makeatletter
\newlabel{lemma-bivariant-zero}{{02UC}{source}{}{sourcepointer-6}{}}
\newlabel{lemma-bivariant-zero@cref}{{[originallemma][0][]02UC}{[1][0][]source}{}{}{}}
\makeatother
\crefname{originallemma}{original lemma}{original lemmas}
\Crefname{originallemma}{Original lemma}{Original lemmas}
\makeatletter
\newlabel{lemma-gysin-fundamental}{{0FBM}{source}{}{sourcepointer-7}{}}
\newlabel{lemma-gysin-fundamental@cref}{{[originallemma][0][]0FBM}{[1][0][]source}{}{}{}}
\makeatother
\title{For a flat local complete intersection morphism the Gysin class is the flat pullback}
\author{The Stacks Project Authors}
\date{stacks-project:chow.tex:local-snapshot-20261012; GNU FDL 1.2 or later}
\begin{document}
\maketitle
\noindent\textit{Editorial note.} All delivered mathematical text is the author's own, copied from the pinned Stacks Project chapter file chow.tex; no statement or proof step was rewritten, reordered or removed.\par\medskip
\noindent\textit{Editorial note.} Preamble provenance: the pinned chapter file opens with an input of the shared preamble and carries no preamble text of its own, so the project's single Stacks-wide preamble staged at sources/stacks-chow/preamble.tex is excerpted verbatim as two spans, lines 14-18 (Xy-pic setup) and lines 154-217 (the native theorem declarations and mathematical macros). Nothing was added, removed or redefined; the wrapper class and the amsmath/amssymb/amsthm/graphicx base packages are supplied by the assembler.\par\medskip
\noindent\textit{Editorial note.} The counted statement and proof contain no cross reference, no citation and no included file, so the delivered body is self contained.\par\medskip
\noindent\textit{Editorial note.} Canonical location of the counted proposition: Stacks tag 0FF4 (sources/stacks-chow/pinned-original-source.tex line 13547).\par\medskip
\noindent\textit{Allowed standard prerequisites (editorial).}\begin{itemize}
\item The flat pullback as a bivariant class. Same chapter, Lemma labelled lemma-flat-pullback-bivariant, invoked in the statement. Supplies the class that f\textasciicircum{}! is shown to equal.
\item Testing identities of bivariant classes on an integral base. Same chapter, Lemma labelled lemma-bivariant-zero, cited at proof lines 13563-13565. Reduces the proof to f\textasciicircum{}*[Y] = i\textasciicircum{}!(g\textasciicircum{}*[Y]) with Y integral.
\item Base change of smooth, syntomic and flat local complete intersection morphisms. Morphisms, Lemmas labelled morphisms-lemma-base-change-smooth and morphisms-lemma-base-change-syntomic, and More on Morphisms, Lemma labelled more-morphisms-lemma-flat-lci, cited at proof lines 13557-13561. Allows the reduction to an integral target.
\item The Gysin class of a regular immersion on the fundamental class. Same chapter, Lemma labelled lemma-gysin-fundamental, cited at proof lines 13570-13572. Gives i\textasciicircum{}![P]\_\{n+t\} = [X]\_\{n+s\}.
\end{itemize}\medskip
\ifdefined\hypertarget\hypertarget{sourcepointer-1}{}\fi
\noindent\textit{Source reference.} Original reference 02QL: sources/stacks-chow/pinned-original-source.tex, source lines 1372--1372. Same-chapter pointer into the author environment situation at source line 1372; invoked at proof line None of the counted proof. It is retained only as a navigation pointer; every load-bearing step is inside the delivered span.\par\medskip
\ifdefined\hypertarget\hypertarget{sourcepointer-2}{}\fi
\noindent\textit{Source reference.} Original lemma 0B78: sources/stacks-chow/pinned-original-source.tex, source lines 6131--6131. Same-chapter pointer into the author environment lemma at source line 6131; invoked at proof line None of the counted proof. It is retained only as a navigation pointer; every load-bearing step is inside the delivered span.\par\medskip
\ifdefined\hypertarget\hypertarget{sourcepointer-3}{}\fi
\noindent\textit{Source reference.} Original lemma 01VB: sources/stacks-chow/tags.txt, source lines 2440--2440. Cross-chapter standard prerequisite located in the project tag list staged in the local snapshot; invoked at proof line 13560 of the counted proof. It is retained only as a navigation pointer; every load-bearing step is inside the delivered span.\par\medskip
\ifdefined\hypertarget\hypertarget{sourcepointer-4}{}\fi
\noindent\textit{Source reference.} Original lemma 01UI: sources/stacks-chow/tags.txt, source lines 2411--2411. Cross-chapter standard prerequisite located in the project tag list staged in the local snapshot; invoked at proof line 13561 of the counted proof. It is retained only as a navigation pointer; every load-bearing step is inside the delivered span.\par\medskip
\ifdefined\hypertarget\hypertarget{sourcepointer-5}{}\fi
\noindent\textit{Source reference.} Original lemma 069K: sources/stacks-chow/tags.txt, source lines 8044--8044. Cross-chapter standard prerequisite located in the project tag list staged in the local snapshot; invoked at proof line 13562 of the counted proof. It is retained only as a navigation pointer; every load-bearing step is inside the delivered span.\par\medskip
\ifdefined\hypertarget\hypertarget{sourcepointer-6}{}\fi
\noindent\textit{Source reference.} Original lemma 02UC: sources/stacks-chow/pinned-original-source.tex, source lines 6585--6585. Same-chapter pointer into the author environment lemma at source line 6585; invoked at proof line 13564 of the counted proof. It is retained only as a navigation pointer; every load-bearing step is inside the delivered span.\par\medskip
\ifdefined\hypertarget\hypertarget{sourcepointer-7}{}\fi
\noindent\textit{Source reference.} Original lemma 0FBM: sources/stacks-chow/pinned-original-source.tex, source lines 12076--12076. Same-chapter pointer into the author environment lemma at source line 12076; invoked at proof line 13570 of the counted proof. It is retained only as a navigation pointer; every load-bearing step is inside the delivered span.\par\medskip
\begin{lemma}
\label{lemma-lci-gysin-flat}
Let $(S, \delta)$ be as in Situation \ref{situation-setup}.
Let $f : X \to Y$ be a local complete intersection morphism
of schemes locally of finite type over $S$. If the Gysin map
exists for $f$ and $f$ is flat, then $f^!$ is equal to the
bivariant class of Lemma \ref{lemma-flat-pullback-bivariant}.
\end{lemma}

\begin{proof}
Choose a factorization $f = g \circ i$ with $i : X \to P$
an immersion and $g : P \to Y$ smooth. Observe that for
any morphism $Y' \to Y$ which is locally of finite type,
the base changes of $f'$, $g'$, $i'$ satisfy the same
assumptions (see Morphisms, Lemmas \ref{morphisms-lemma-base-change-smooth}
and \ref{morphisms-lemma-base-change-syntomic} and
More on Morphisms, Lemma \ref{more-morphisms-lemma-flat-lci}).
Thus we reduce to proving that $f^*[Y] = i^!(g^*[Y])$ in case $Y$
is integral, see Lemma \ref{lemma-bivariant-zero}. Set $n = \dim_\delta(Y)$.
After decomposing $X$ and $P$ into connected components we
may assume $f$ is flat of relative dimension $r$ and
$g$ is smooth of relative dimension $t$.
Then $f^*[Y] = [X]_{n + s}$ and $g^*[Y] = [P]_{n + t}$.
On the other hand $i$ is a regular immersion of codimension $t - s$.
Thus $i^![P]_{n + t} = [X]_{n + s}$ (Lemma \ref{lemma-gysin-fundamental})
and the proof is complete.
\end{proof}

\end{document}
